\subsection{Что результаты означают для архитектуры}
Архитектурный результат касается размещения предписанных функций. SR объединяет инструкции планирования, реализации и проверки в одном вызове и требует меньше наблюдаемых ресурсов, чем его трёхвызовный вариант MR. Контраст SR--SN отвечает на другой вопрос: меняются ли обозначения ролей при одинаковых операциях и числе вызовов. Меньший средний расход не доказывает улучшение качества, связанное с ролевыми обозначениями. Более того, D использует меньше ресурсов, чем SR, внутри основной факторной серии. Настоящее исследование оценивает эти решения как компоненты ядра Hybrid-INoT; данные позволяют рассматривать компактную конфигурацию одного вызова, оставляя пользу ролевой разметки недоказанной.

Внешняя проверка нужна для интерпретации исхода: она отличает сгенерированную программу от программы, прошедшей заданные исполняемые тесты. Поскольку оценщик одинаков для условий, его наличие не объясняет различие, специфичное для Hybrid-INoT. Независимая оценка тестами также не доказывает, что система исправит неуспех без дополнительных условий. Отдельное сравнение INoT* расширяет архитектурный вопрос на другую инструкцию внутреннего диалога, не меняя определения ядра SR.

Завершённое сравнение SCC позволяет оценить ещё одну внешнюю процедуру. Её денежная оценка по тарифам API примерно в 3{,}14 раза выше, чем у SR, и в 2{,}94 раза выше, чем у SN; здесь сравниваются средние затраты по парным задачам. Качество на наблюдаемых парах ниже на 2{,}68 и 2{,}48 процентного пункта соответственно. Эти условные разности проходят поправку Холма для двух тестов, но предельные оценки с учётом пропусков допускают любой знак. Вывод относится к данной конфигурации авторского кода SCC и наблюдаемой выборке; порядок качества во всей назначенной выборке не установлен.

\subsection{Расширения дизайна и необходимые доказательства}
\begin{table}[htbp]\centering\small
\caption{Архитектурные вопросы и границы имеющихся данных. Это структура обсуждения, а не дополнительные подтверждающие тесты.}
\label{tab:hybrid-scope}
\begin{tabularx}{\textwidth}{@{}>{\raggedright\arraybackslash}p{0.23\textwidth}>{\raggedright\arraybackslash}X>{\raggedright\arraybackslash}X@{}}\toprule
Вопрос & Текущий вклад & Необходимые данные\\\midrule
Эффективность общего контекста с сохранением качества & Условный баланс токенов; сравнения качества и ресурсов компонентов без компрессии & Контролируемое варьирование длины контекста и внешне обоснованный критерий сохранения качества\\
Экономическая ценность предварительного отбора & Перспективное условие безубыточности & Интегрированная политика отбора и повторов с итоговыми проверенными исходами и всеми затратами\\
Граница при параллельных инструментах & Явно вне фиксированной текстовой генерации & Сопоставимые последовательные и параллельные инструменты на независимых подзадачах; задержка и проверенное качество\\\bottomrule
\end{tabularx}\end{table}

Например, пусть малая предварительная проверка стоит $c_s$, дорогой повтор имеет постоянную стоимость $c_r$, а вероятности повтора при двух политиках равны $\rho_0$ и $\rho_1$. При ровно одной предварительной проверке и не более чем одном повторе изменение ожидаемой стоимости генерации составляет
\begin{equation}
\Delta C=c_s+(\rho_1-\rho_0)c_r.
\label{eq:hybrid-screen}
\end{equation}
Расширенная политика дешевле тогда и только тогда, когда $c_s<(\rho_0-\rho_1)c_r$. Это условное тождество сохраняет экономический проектный вопрос, но не оценивает эффект политики: проверяющая модель может сократить повторы, просто принимая неверные кандидаты. Полезная политика требует также проверенного качества и заранее обоснованного правила принятия. В основной серии ни малая проверяющая модель, ни политика повторов не выполняются.

Аналогично независимые инструменты могут выигрывать от внешнего параллелизма, даже если единый ролевой вызов использует меньше токенов. Текущая последовательность этапов не проверяет этот вопрос задержки. Компрессия, адаптивный выбор режима и коррекция по результатам проверки изменили бы воздействие и потребовали бы отдельных сравнений. Это возможные расширения Hybrid-INoT, а не проверенные компоненты, скрытые внутри сообщаемой ресурсной разности.
